\jabstract{
近年我が国では、多くの地域において少子高齢化や人口流出により労働人口が減少し、地域課題の解決が困難になりつつある。これに対し、文部科学省は地方大学にその解決への取り組みを期待している面がある。実際、多くの大学も自身の存在意義を示すために地域に根ざした研究に力を入れている。一方で、計量書誌学の観点からは、このような地域研究の活動を示すような指標の開発はあまり見られない。これにはいくつかの問題があるからだと考えられる。まず、地域の名称（地名）を文章等から正しく抽出することが困難であり、つぎにその地名が唯一に決められるかも難しい問題である。最終的には対象となる論文やデータに対して地名や地理情報をユニークに紐づける（キュレーションする）仕組みが必要となる。本研究では、このようなソリューションを視野に入れるものの、まずは地名の抽出という問題に着目し、コンパクトなテストデータを用いた試みを行う。
}
